% TC-0013 test-case definition (WSP test-case template).
\documentclass[letterpaper]{article}

\usepackage[letterpaper,margin=1in]{geometry}

% Compile from the repository root so this path resolves to the pinned WSP
% submodule.
\input{wsp/testing/test-case-library.tex}

\begin{document}

\def\TCID{TC-0013}
\def\TCTitle{Corpus 218: MonthView GetDayBold event}
\def\TCPurpose{%
  Verify that GetDayBold is raised for the stretch of days the calendar shows, that State is an array from 0 to Count - 1, that the days the handler sets True are bold, and that the others are not (even a day made plain in Form\_Load).}
\def\TCPriority{\TCPriorityHigh}

\def\TCRequirementRef{REQ-0003}
\def\TCDesignRef{Not applicable.}
\def\TCFeatureRef{\texttt{tests/corpus/218-mscomctl2-getdaybold.vbp}, \texttt{.frm}, \texttt{.ui}, \texttt{.debug} and \texttt{.expected}.}
\def\TCTestTechnique{%
  Scenario test: a program that prints what it observes, compared with a stored transcript.}

\def\TCPreconditions{%
  A \texttt{wfc} executable that includes this library (WFC version 0.8.28 or later) and the files of the case in \texttt{tests/corpus}.}
\def\TCEnvironment{%
  Windows x64. The \texttt{wfc} executable of WFC, given to CMake as \texttt{WFC\_EXE}; the program is run by the interpreter (\texttt{--interpret}) with the system common controls. The case is registered only when \texttt{WFC\_EXE} is set; without it CTest reports the case as skipped.}
\def\TCAssumptions{%
  Not applicable.}

\def\TCInputData{%
  The script line \texttt{ui Form1.Month1 getdaybold 2027-06-01 30}; DayBold(2027-06-06) set to False in Form\_Load; the days 2027-06-06, 06-07, 06-13, 06-27 and 07-04 read in Form\_Unload.}
\def\TCInitialState{%
  No output of an earlier run is used. CTest test name: \texttt{TC-0013-corpus-getdaybold}; when another project mounts the library the name is prefixed with \texttt{wfc-mscomctl2-}.}

\def\TCProcedure{%
  \begin{enumerate}
    \item Run \texttt{wfc --interpret tests/corpus/218-mscomctl2-getdaybold.vbp}. The handler sets State True for every Sunday of the stretch and prints what it was asked.
    \item The script asks for 30 days from 2027-06-01 and closes the form; Form\_Unload prints DayBold of five days.
    \item Compare the Immediate-window output with \texttt{218-mscomctl2-getdaybold.debug} and standard output with the empty \texttt{218-mscomctl2-getdaybold.expected}.
  \end{enumerate}}
\def\TCExpectedResult{%
  The first line is \texttt{ask2027-06-01 30  0  29} (with the spacing of Debug.Print); the bold line is \texttt{boldTrueFalseTrueTrueFalse}. Both outputs equal the stored files.}
\def\TCPostConditions{%
  No files are written; the interpreter's windows are closed by the script.}

\TCTemplate

\end{document}
